\section{Unified Treatment of Three Feedback Loops}
\label{sec:unified_loops}

The three compound feedback loops identified in the gap analysis across
the prior papers in this sequence~\citep{lasser2026gfm,lasser2026poset,lasser2026horizon,lasser2026scm,lasser2026exo} are instances of the same $(P, \Wm)$ dynamical system with
different effective parameters.  This section maps each loop to the
formalism and shows how the phase boundary specializes.

\subsection{Loop 1: Structural avoidance $\to$ leverage blindness}

\paragraph{Mapping.}
The ``subsumption'' here is not agent elimination but \emph{structural
discovery avoidance}: the actor avoids learning poset structure
(subsumption edges) because discovery can contract $\volL$ in the short
term.

\begin{itemize}
\item $P_t$: the actor's \emph{known} poset (a subset of the true
  poset).
\item $\Wm_t$: leverage estimates $\hat{\lambda}(d, t)$, which are
  accurate only for the known poset.
\item ``Observation channels'' lost: each avoided structural discovery
  removes a potential correction to $\hat{\lambda}$.
\item $\rW$: the rate at which leverage estimates drift due to
  un-learned structure.
\item $\rS$: the rate at which the structural discovery
  value~\citep[Proposition~2, Definition~5]{lasser2026horizon} drives the actor toward
  learning structure.
\end{itemize}

\paragraph{Phase boundary specialization.}
Self-correction dominates when the information value $I(d_1 \leq d_2)$
\citep[Definition~5]{lasser2026horizon}) exceeds the immediate $\volL$ contraction from
discovery for a sufficient fraction of unknown edges.  The critical
ratio becomes:
\begin{equation}
\label{eq:loop1_boundary}
\min_{\text{undiscovered edges}}
  \frac{I(\text{edge})}{|\Delta\volL(\text{discover edge})|}
  > \text{threshold}
\end{equation}
where the threshold is determined by the degradation and drift
parameters from the general phase boundary.  This formalizes the
discussion of structural avoidance in~\citet{lasser2026poset} as a phase-boundary condition
rather than an informal concern.

\subsection{Loop 2: Partial subsumption $\to$ world-model homogenization}

This is the primary case, directly formalized in
Sections~\ref{sec:coupled_system}--\ref{sec:phase_boundary}.  The
parameters are as defined throughout: $P_t$ is the capability poset,
$\Wm_t$ is the full world-model accuracy vector, observation channels
are rooted in agents, and the phase boundary is stated in terms of
$(\rS, \rW, \rsub, \dmax, \rhomincross)$.

\subsection{Loop 3: Over-calibrated restriction $\to$ discovery avoidance}

\paragraph{Mapping.}
The ``subsumption'' is the actor \emph{restricting} capabilities rather
than eliminating agents.  But the dynamical structure is identical:
restriction prevents exercise, which prevents observation, which
prevents calibration of the risk estimates that justified the
restriction.

\begin{itemize}
\item $P_t$: the set of unrestricted capabilities (restriction reduces
  the exercisable poset).
\item $\Wm_t$: risk estimates $\hat{\Risk}(S, t)$, which are accurate
  only for exercised capabilities.
\item ``Observation channels'' lost: each restriction prevents the
  exercise events that would have generated Bayesian evidence about
  the risk claim (the verification asymmetry from
  \citet{lasser2026horizon}, Remark~2).
\item $\rW$: the rate at which risk estimates drift due to
  un-exercised pathways (this is exactly the Wamura problem from
  \citet{lasser2026horizon}).
\item $\rS$: the rate at which controlled relaxation protocols or
  natural exercise events provide counter-evidence.
\end{itemize}

\paragraph{Phase boundary specialization.}
Self-correction dominates when the controlled relaxation protocol
generates evidence faster than restriction accumulates unfalsifiable
risk claims.  This connects the present paper to the Wamura problem: a
controlled relaxation protocol is the \emph{mechanism} that makes $\rS$
positive for Loop~3, without which $\rS = 0$ and the absorbing basin is
the entire state space.

\begin{proposition}[Wamura Problem as Phase Boundary Degeneracy]
\label{prop:wamura}
In the absence of controlled relaxation ($\rS = 0$ for Loop~3), the
phase boundary collapses: every non-trivial risk restriction is
absorbing, and the framework converges to maximal restriction.
\end{proposition}

\begin{proof}
With $\rS = 0$, the cumulative error budget $B(t_0, T)$ becomes
\begin{equation*}
B(t_0, T) = \sum_{t \in \mathrm{sub}} \dmax \cdot \Delta\rho(t)
  + \sum_{t \in \mathrm{sub}} \rW \cdot \Lyap(\Wm_t)
\end{equation*}
Both sums contain only non-negative terms ($\dmax, \Delta\rho, \rW,
\Lyap \geq 0$), and the correction term vanishes entirely.
Therefore $B(t_0, T) \geq 0$ for every window, with equality only in
the trivial case where no restriction occurs.

For any non-trivial restriction step ($\nsub(t) = 1$ with
$\Delta\rho(t) > 0$ or $\Lyap(\Wm_t) > 0$), $B > 0$, satisfying
condition~(i) of Theorem~\ref{thm:phase_boundary}, Part~(b).
Conditions~(ii) and~(iii) follow from the restriction-specific
dynamics: each restriction step removes the exercise that would
generate evidence (no exercise $\to$ no evidence), driving
$\rhocross_k \to 0$ on the risk dimension (satisfying~(ii)).  With
$\rS = 0$ and the risk dimension environmentally non-stationary (the
true risk changes as the restricted capability's context evolves
without observation), the error $\epsilon_k$ grows, eventually
exceeding $\epsilon_{\mathrm{crit},k}$ (satisfying~(iii)).  Note that
$\dmax > 0$ bounds the drift rate from above (C2$'$); the cascade
requires positive drift, which is a structural consequence of
non-stationarity on the restricted dimension.  All three conditions of Part~(b) are
therefore met, and the cascade drives the system toward maximal
restriction.

This is the formal statement of why the Wamura problem was called
``the most important open problem'' in the discussion of~\citet{lasser2026horizon}: without a
mechanism to generate evidence against unfalsifiable risk claims, the
feedback loop between restriction and evidence suppression is
absorbing by default.
\end{proof}

\subsection{Summary}

Table~\ref{tab:loops} summarizes the three loops as parameter
instantiations of the unified system.

\begin{table}[h]
\centering
\small
\begin{tabular}{@{}lllll@{}}
\toprule
\textbf{Loop} & \textbf{$P_t$} & \textbf{$\Wm_t$} &
\textbf{$\rS$ source} & \textbf{$\rW$ source} \\
\midrule
1. Structural avoidance & Known poset & Leverage estimates &
  Discovery value (P3 Def.\ 5 / Prop.\ 2) & Un-learned structure \\
2. Partial subsumption & Capability poset & Full accuracy vector &
  $T_s$ self-trust (P1) & Channel loss \\
3. Over-calibrated restriction & Unrestricted caps. & Risk estimates &
  Controlled relaxation & Unfalsifiable risk \\
\bottomrule
\end{tabular}
\caption{Three feedback loops as instances of the $(P, \Wm)$ system.}
\label{tab:loops}
\end{table}

The unified treatment reveals that solutions for one loop improve the
phase boundary for all three.  Controlled relaxation (which makes
$\rS > 0$ for Loop~3) has no direct effect on Loop~2's parameters, but
the structural insight transfers: any mechanism that generates
observation-channel signal on dimensions that would otherwise go blind
is a contribution to $\rS$ in the relevant loop's phase boundary.
